% #############################################################################
% This is Appendix B
% !TEX root = ../main.tex
% #############################################################################
\chapter{LoRA Rank Ablation}
\label{chap:appendix-lora-rank}

This appendix reports a LoRA-rank ablation for the decoder-controlled
4B-parameter T5Gemma 2 model. The comparison examines how performance changes
across LoRA ranks under the supervised training setup used for the corresponding
Stage B GPT-Wikipedia configuration.

\section{Ablation Design}

The 4B-parameter experiments use LoRA rank 48. This ablation compares ranks 8,
16, 32, and 48 in the decoder-controlled configuration trained with Stage A on
OpenSubtitles followed by Stage B on GPT-Wikipedia. All runs use the supervised
protocol selected in \Cref{sec:smaller-model-protocol-study}. The control
strings, data splits, target modules, learning rates, dropout, and evaluation
protocol are kept fixed, together with the LoRA scaling factor $\alpha = 96$.

{\tolerance=1500\looseness=-1
However, the implementation scales the LoRA update by $\alpha/r$, so holding
$\alpha$ fixed means that the effective adapter multiplier changes with rank.
The resulting multipliers are 12, 6, 3, and 2 for ranks 8, 16, 32, and
48, respectively. The comparison therefore varies rank together with this
effective scaling and should not be interpreted as an isolated effect of rank
alone.
\Cref{tab:appendix-lora-rank} reports the results on the Golden Collection
dataset, as it is the benchmark most closely aligned with the conservative
rewriting objective.\par}

\section{Results}

\begin{table}[htbp]
\centering
\thesistablefont
\setlength{\tabcolsep}{5pt}
\caption{Decoder-controlled LoRA rank ablation on the Golden Collection.}
\label{tab:appendix-lora-rank}
\begin{tabular}{cccccc}
\toprule
\multirow{2}{*}{\textbf{Rank}} & \multicolumn{2}{c}{\textbf{Trainable parameters}} & \multirow{2}{*}{\shortstack{\textbf{Macro-F1}\\\textbf{(\%)}}} & \multirow{2}{*}{\shortstack{\textbf{BLEU}\\\textbf{(0--100)}}} & \multirow{2}{*}{\textbf{TER}} \\
\cmidrule(lr){2-3}
& \textbf{Count} & \textbf{\% of rank 48} & & & \\
\midrule
8  & 31.3M  & 16.7\%  & 70.95 & 87.14 & 0.0943 \\
16 & 62.6M  & 33.3\%  & \textbf{72.31} & 87.24 & 0.0964 \\
32 & 125.2M & 66.7\%  & 71.17 & 87.23 & 0.0944 \\
48 & 187.8M & 100.0\% & 71.53 & \textbf{87.68} & \textbf{0.0895} \\
\bottomrule
\end{tabular}
\end{table}

Rank 48 gives the best observed rewriting performance, with the highest BLEU
(87.68) and lowest TER (0.0895). However, its BLEU advantage is modest: 0.44
points over rank 16, the next-highest result, and 0.54 points over rank 8.
Ranks 8--32 are tightly clustered between 87.14 and 87.24 BLEU, while rank 8
uses approximately one-sixth as many trainable adapter parameters as rank 48.
Identification does not follow the same ordering: rank 16 gives the highest
macro-F1, at 72.31 compared with 71.53 for rank 48, and macro-F1 does not
increase monotonically with rank. The results therefore show a trade-off
between trainable adapter size and performance. Rank 48 performs best on the
rewriting metrics in this configuration, while substantially smaller adapters
remain competitive when trainable-parameter efficiency is prioritized.

The experiment does not establish statistical significance or a general
monotonic relationship between LoRA rank and performance.

{\tolerance=1500\looseness=-1
Reducing LoRA rank lowers the number of trainable adapter parameters, but it
does not reduce the size of the frozen base model or proportionally reduce
activation and forward-pass costs. The percentages in
\Cref{tab:appendix-lora-rank} describe adapter-parameter savings, not
proportional reductions in total \mbox{training memory or computation}.
\par}
